\chapter{Conclusão}
\label{cap:conclusao}

Este trabalho investigou o uso de Modelos de Linguagem de Grande Porte para simplificação e adaptação comunicacional de textos técnicos de saúde em um contexto mHealth. A motivação central partiu da dificuldade que conteúdos de saúde, frequentemente escritos em linguagem técnica, podem impor a usuários finais de aplicações digitais, especialmente quando esses conteúdos envolvem orientações de cuidado que precisam ser compreendidas de forma clara. Nesse cenário, o objetivo geral do trabalho foi investigar, implementar e avaliar uma solução baseada em LLMs para simplificação e adaptação de linguagem técnica em uma plataforma mHealth baseada em planos de cuidado, utilizando guardrails para apoiar a preservação do significado e das informações essenciais do material original. Esse objetivo foi desenvolvido a partir da fundamentação teórica, da análise de trabalhos relacionados, da proposta da solução, da implementação funcional da aplicação SimplyHealth e da avaliação com materiais de saúde previamente definidos, perfis simulados, indicador de legibilidade, análise qualitativa e verificação da atuação dos guardrails.

A proposta buscou enfrentar a dificuldade de compreensão de conteúdos técnicos de saúde por meio de um fluxo de simplificação textual orientado por perfis comunicacionais de pacientes. Nesse contexto, a Takere foi utilizada como ambiente de referência para situar a abordagem em um cenário mHealth baseado em planos de cuidado, sem que tenha sido realizada uma integração plena com a plataforma. A adaptação foi delimitada como comunicacional, e não clínica: os perfis simulados, definidos a partir de aspectos como idade, escolaridade e área de formação, foram utilizados apenas para ajustar vocabulário, tom, nível de explicação e forma de apresentação do conteúdo. Dessa forma, a LLM foi empregada como mecanismo de reescrita e adaptação linguística, não como fonte autônoma de conhecimento médico, nem como agente responsável por validar, recomendar ou modificar condutas clínicas.

Essa proposta foi materializada em uma implementação funcional, com frontend e backend, capaz de receber textos livres ou arquivos nos formatos PDF e TXT, processar o conteúdo original e gerar saídas voltadas à compreensão do usuário. A aplicação implementou a geração de texto simplificado, glossário, perguntas frequentes, chat contextual e apresentação do indicador de legibilidade Flesch-PT. Além disso, os prompts foram estruturados para orientar a tarefa da LLM e incorporar regras de preservação do conteúdo, enquanto os guardrails foram operacionalizados por meio de instruções no prompt, verificações determinísticas, avaliação por LLM e mecanismo de fallback. Essa implementação demonstra que o trabalho não se limitou a uma proposta conceitual, mas resultou em uma solução funcional voltada à investigação prática do uso de LLMs na simplificação de linguagem técnica em saúde.

A avaliação da solução considerou dois materiais de saúde previamente definidos: um sobre estomia e outro sobre cuidados com os pés para pessoas com diabetes. Foram utilizados três perfis simulados, Ana, Carlos e Maria, representando diferentes necessidades comunicacionais, e realizadas 18 execuções principais. Os resultados de legibilidade indicaram aumento do Flesch-PT em 15 dos 18 casos avaliados. A média dos textos originais foi de 14,5 pontos, enquanto a média das versões simplificadas foi de 24,4 pontos, resultando em um ganho médio geral de +9,9 pontos. Maria apresentou os maiores ganhos de legibilidade, o que se mostrou coerente com seu perfil de menor escolaridade e maior necessidade de linguagem acessível. Carlos apresentou ganhos intermediários, enquanto Ana apresentou os menores ganhos e, no material sobre diabetes, variações negativas, interpretadas com cautela como possível consequência da preservação de linguagem mais técnica para um perfil com formação na área da saúde.

Além dos resultados quantitativos de legibilidade, a avaliação também evidenciou o papel ativo dos guardrails na filtragem de saídas problemáticas. Em diferentes tentativas, os mecanismos de verificação rejeitaram alterações consideradas inadequadas, como perda de negação, troca de substância, alteração de instrução procedimental, contradição numérica e mudança de sentido. Esses casos reforçam a importância de não tratar a simplificação textual em saúde apenas como uma tarefa de tornar frases mais fáceis, mas como um processo que exige controle sobre a preservação das informações essenciais do material original. Assim, os resultados indicam a viabilidade da abordagem no contexto avaliado, tanto pela capacidade de gerar versões mais acessíveis em parte significativa dos casos quanto pela atuação dos mecanismos de controle sobre alterações indevidas.

Com relação à pergunta de pesquisa - como Modelos de Linguagem de Grande Porte podem ser utilizados, com apoio de guardrails, para simplificar e adaptar textos técnicos de saúde ao perfil de pacientes, preservando a precisão médica das informações apresentadas - os achados sugerem que essa utilização é viável quando a geração é baseada em textos previamente definidos, o perfil do paciente é usado apenas para adaptação comunicacional e há restrições explícitas sobre o que a LLM pode ou não modificar. Dentro do escopo avaliado, a solução demonstrou que LLMs podem apoiar a reescrita de conteúdos técnicos de saúde de modo sensível a diferentes perfis de usuários, desde que combinadas com prompts bem delimitados, guardrails de verificação e análise das informações essenciais preservadas. No contexto deste trabalho, portanto, a preservação da precisão médica deve ser compreendida como a manutenção do sentido e das orientações presentes no material original, e não como validação clínica autônoma do conteúdo. Dessa forma, não é possível afirmar que a solução garante precisão médica, que os textos foram clinicamente validados, que pacientes reais compreenderiam melhor os conteúdos ou que a abordagem melhora adesão ao tratamento ou desfechos de saúde.

As principais contribuições deste trabalho estão na proposição, implementação e avaliação de uma solução funcional para simplificação e adaptação comunicacional de textos técnicos de saúde com LLMs. A SimplyHealth demonstrou um fluxo em que perfis simulados orientam diferenças de linguagem, vocabulário, tom e nível de explicação, mantendo a adaptação restrita à forma de comunicação. Outra contribuição está na aplicação de guardrails para reduzir riscos de alterações indevidas no conteúdo original, principalmente em um domínio no qual pequenas mudanças de sentido podem ter consequências relevantes. O trabalho também contribui ao integrar, em uma mesma solução, texto simplificado, glossário, perguntas frequentes, chat contextual e indicador de legibilidade, permitindo observar a simplificação não apenas como uma saída textual isolada, mas como parte de uma experiência de apoio à compreensão em mHealth. Por fim, a avaliação funcional e linguística, baseada no Flesch-PT, em exemplos qualitativos e na análise da atuação dos guardrails, permitiu discutir benefícios e limites da abordagem de forma prática e contextualizada.

Apesar dessas contribuições, os resultados devem ser interpretados considerando as limitações do estudo. A avaliação utilizou apenas dois materiais de saúde e três perfis simulados, o que restringe a generalização dos achados para outros tipos de conteúdo, condições de saúde e públicos. Além disso, não houve avaliação com pacientes reais nem validação clínica formal por especialistas, de modo que não é possível afirmar que as versões simplificadas seriam compreendidas adequadamente em situações reais de cuidado. O Flesch-PT, embora útil como indicador automático de legibilidade, mede apenas aspectos formais do texto e não substitui avaliações de compreensão, clareza percebida ou adequação clínica. Também é importante reconhecer que os guardrails reduzem riscos, mas não garantem preservação absoluta do conteúdo original, e que os resultados dependem do modelo utilizado, dos prompts, dos materiais avaliados e do ambiente de execução. A ausência de integração plena com a Takere também delimita o alcance da solução no contexto de uma plataforma mHealth real.

Como trabalhos futuros, recomenda-se ampliar a avaliação para um conjunto maior de materiais de saúde, incluindo conteúdos com diferentes níveis de complexidade e criticidade. Também seria relevante realizar estudos com pacientes reais ou usuários finais, investigando compreensão textual, percepção de clareza e dificuldades de uso em situações mais próximas do contexto real. Outra possibilidade é envolver profissionais de saúde na revisão sistemática das saídas geradas, permitindo uma análise mais robusta da preservação das informações clínicas. Além disso, futuras versões da solução poderiam comparar diferentes modelos de linguagem, aprimorar os guardrails determinísticos e baseados em LLM, melhorar mecanismos de rastreabilidade das simplificações e avaliar uma integração mais próxima com a Takere. Assim, este trabalho contribui ao demonstrar, por meio de uma implementação funcional, que LLMs podem apoiar a adaptação comunicacional de textos técnicos de saúde, desde que utilizadas com escopo delimitado, mecanismos de controle e avaliação cuidadosa das informações preservadas.
